\section*{Введение}
\addcontentsline{toc}{section}{Введение}

\textbf{Актуальность темы.} Феномен политической социализации играет фундаментальную роль в обеспечении стабильности, воспроизводства и адаптации любой политической системы. Институциональный каркас государства, его правовые нормы и управленческие алгоритмы не способны эффективно функционировать без соответствующего укоренения в ценностных ориентациях, убеждениях и поведенческих установках граждан. В условиях динамичных глобальных трансформаций, кризиса традиционных идеологий и стремительного развития цифровых коммуникаций механизмы интеграции индивида в политическое пространство претерпевают качественные изменения. 

Для современной России данная проблематика приобретает особую остроту. Пережив в конце XX века радикальный слом советской государственности и нормативно-ценностной системы, российское общество столкнулось с затяжным периодом ценностной аномии и фрагментации политической культуры. Эволюция институтов социализации в современной России происходит на фоне геополитического противостояния, необходимости укрепления государственного суверенитета и формирования новых ориентиров общенациональной идентичности. Исследование специфики, кризисов и ключевых агентов политической социализации в российском контексте позволяет эксплицировать латентные пружины электорального поведения, динамику молодежного активизма и векторы легитимации публичной власти, что имеет критическое значение для прогнозирования устойчивости макрополитической системы.

\textbf{Объект исследования} --- политическая социализация как непрерывный процесс интеграции индивида в политическую систему и трансляции ценностно-нормативного опыта между поколениями.

\textbf{Предмет исследования} --- структурно-функциональные особенности, институциональные агенты и ключевые тенденции эволюции процессов политической социализации в современном российском обществе.

\textbf{Цель работы} --- осуществить комплексный теоретико-методологический и прикладной политологический анализ механизмов, кризисных факторов и специфики трансформации политической социализации в условиях современной российской реальности.

Для достижения поставленной цели необходимо решить следующие \textbf{задачи}:
\begin{enumerate}
    \item Исследовать понятийный аппарат, структуру, этапы и базовые типы политической социализации в классической и современной политической науке.
    \item Описать систему традиционных и новейших агентов политической социализации, оценив их функциональную специфику.
    \item Проанализировать исторический контекст, преемственность ценностей и детерминанты формирования нормативных матриц в российском транзитном социуме.
    \item Выявить ключевые кризисные макрофакторы, институциональные деформации и тенденции цифровизации процессов социализации в современной России.
\end{enumerate}